% mage-auto-heading: modeling-behavior-with-states-and-transitions
% mage-auto-heading: ownership-under-contention
% mage-auto-heading: one-representation-several-questions
% Authoritative LaTeX manuscript; Markdown is a generated compatibility projection.

% <!-- part-title: Modeling -->

% <!-- chapter-title: Behavioral Models: What May Happen? -->

\chapter{Behavioral Models: What May Happen?}

A behavioral model represents how a system acts or changes over time. It may describe the messages exchanged by components, the order of steps in a workflow, or how a quantity changes continuously. Where a structural model like a dependency graph answers \emph{what is connected to what?}, a behavioral model asks \emph{what may happen, and in what order?} The representation should preserve the distinctions needed to answer the engineering question while leaving out implementation details that do not affect it.

\section{Modeling Behavior with States and Transitions}\label{heading:modeling-behavior-with-states-and-transitions}

% <!-- point: behavior-has-a-long-repertoire-of-state-transition-representations | State-transition models make lifecycle and recovery questions explicit; state machines, labeled transition systems, and statecharts represent behavior, while temporal specifications state properties over executions. -->

For lifecycle and recovery questions, a state-transition model makes explicit which states matter and how the system can move between them. This is the form of behavioral modeling developed here. A finite-state machine names a finite set of states and the transitions between them. A labeled transition system associates transitions with actions or events. Statecharts add hierarchy and concurrency to organize more complex behavior~\citep{harel1987}. Temporal specifications let us state properties over executions, such as whether every started job eventually finishes. These representations support reasoning about order, reachability, and progress without retaining the code that implements each operation.

% <!-- point: the-reduction-is-what-makes-behavior-checkable | Reducing a service to states and transitions turns "can it reach a state it should never occupy?" from a reconstruction over scattered handlers into a question a person or tool can answer. -->

That reduction is what makes behavior checkable. A job-processing service described only through its handlers, database writes, and retries hides the very question an engineer most needs to ask --- can the system reach a state it should never occupy? The same service, reduced to a handful of states and the transitions between them, turns that question into one a person or a tool can actually answer (\autoref{canonical-behavioral}).

\begin{figure}[htbp]
\centering
\includegraphics[width=\linewidth]{book/latex/assets/canonical-behavioral-c9de560202.pdf}
\caption{\emph{The behavioral move: states and transitions.} This example specifies an allowed job lifecycle: a job passes through Running before Succeeded, with retry and failure transitions describing other possible paths. The legend lists example properties to investigate.}
\label{canonical-behavioral}
\end{figure}

% <!-- point: semantics-are-states-and-permitted-transitions-and-the-omissions-are-the-obligations | A behavioral model represents states and transitions; it can specify allowed behavior or describe behavior to be checked against a requirement. Safety and liveness properties ask different questions over the admitted executions. -->

Nodes represent states, and directed edges represent transitions. A path from an initial state describes an execution admitted by the model. The model's purpose tells us how to interpret that execution. A specification declares allowed behavior, as \autoref{canonical-behavioral} does: omitting a direct \texttt{Pending~→~Succeeded} transition forbids that step. A model of an implementation can instead represent behavior we intend to investigate, including states or transitions that may violate a requirement. Representing a state therefore does not, by itself, declare it acceptable.

The required property tells us what to check over those executions. ``A job never reaches \texttt{Succeeded} before \texttt{Running}'' is a \textbf{safety} property: any violation has a finite execution prefix that demonstrates it. ``Every started job eventually finishes or fails'' is a \textbf{liveness} property: observing a finite period of waiting does not establish that the wait will last forever. Reachability asks whether a state can occur; it serves a safety check when the state being sought violates a required invariant. The choice among transition validation, reachability search, invariant checking, and temporal model checking follows from the property we want to establish~\citep{baier-katoen2008}.

A simple job lifecycle also shows the limit of an ordering property. Requiring every path into \texttt{Succeeded} to pass through \texttt{Running} excludes premature completion, but a job could remain in \texttt{Running} indefinitely. These are different questions about the same behavioral model. A transition check asks whether a particular state change is permitted; reachability analysis asks whether some sequence of transitions can lead to an undesirable state; and a temporal property can require that every started job eventually finishes or fails. Investigating that last property requires representing the ways processing can wait or retry, together with assumptions under which progress is expected. The model may therefore need to change as the engineering question changes.

\section{Ownership Under Contention}\label{heading:ownership-under-contention}

% <!-- point: ownership-has-a-long-vocabulary-of-claims-and-leases | Concurrent workers can acquire, lose, and regain rights to a job; locks, leases, capabilities, and claims express different aspects of access and ownership, which the model selects according to the question. -->

When several workers can process the same job, \emph{who owns this job?} depends on their claims. A worker may stop responding, its claim may expire, and another worker may acquire the job. The original worker can later return and attempt to act using its old claim. An ownership model must retain enough information to distinguish the current right to act from a claim that is no longer valid.

\autoref{ownership-constructs} distinguishes four familiar constructs used to represent access and ownership.

{\footnotesize
\begin{longtable}[]{@{}
  >{\raggedright\arraybackslash}p{(\linewidth - 2\tabcolsep) * \real{0.5000}}
  >{\raggedright\arraybackslash}p{(\linewidth - 2\tabcolsep) * \real{0.5000}}@{}}
\caption[Constructs for access and ownership]{\emph{Constructs for access and ownership.} Locks, leases, capabilities, and claims express different aspects of a worker\textquotesingle s right to act.}\label{ownership-constructs}\tabularnewline
\toprule\noalign{}
\begin{minipage}[b]{\linewidth}\raggedright
Construct
\end{minipage} & \begin{minipage}[b]{\linewidth}\raggedright
What it expresses
\end{minipage} \\
\midrule\noalign{}
\endhead
\bottomrule\noalign{}
\endlastfoot
Lock & Coordinates access to a resource. An exclusive lock admits one holder at a time. \\
Lease & Grants a right for a limited period; expiry can make that right available for reassignment~\citep{graycheriton1989}. \\
Capability & Carries authority to perform specified operations on a resource; that authority need not be exclusive. \\
Claim & Records a worker's assignment to a job and the conditions under which that assignment is valid. \\
\end{longtable}
}

These constructs can be combined. For a leased job, the ownership model records the worker's claim and how expiry or reassignment changes it. The next distinction is between successive claims on the same job.

% <!-- point: a-claim-carries-the-facts-that-make-contention-decidable | A claim binds a holder to a resource and carries which holder, which generation of claim, and when it expires — the generation being what tells a current claim from a superseded one — and because the hard failures live in interleavings, the characteristic analysis is state-space exploration with invariant checking. -->

A claim binds a holder to a resource and carries the facts that make contention decidable: which holder, which generation of claim, and when it expires. The generation is the fact a bare "A owns X" edge cannot express --- it distinguishes a \emph{current} claim from a \emph{superseded} one. The difficult failures arise from the interleavings: the orders in which concurrent claims, expiries, and reclaims can occur. So the characteristic analysis is state-space exploration over those sequences, with an invariant checked at each reachable configuration --- the behavioral toolkit of this section pointed at a concurrency question.

For example, a worker may hold a job, lose its claim after a timeout, and later return with an old completion request. The current owner alone is insufficient to decide whether that request still applies: the model must distinguish successive claims on the same job. Recording a claim generation supplies that distinction. At this level, a job is leased exactly when it has a current claim; the claim represents ownership directly. DocAble\textquotesingle s recovery and ownership models show how such small additions support two concrete crash-safety questions.

% <!-- box-family: inset -->

% <!-- inset-domain: WORKED EXAMPLE -->

% <!-- box-weight: deep-dive -->

% mage-box-prefix:

\begin{bookbox}{inset}{Worked example --- Recovery and ownership of a document job}
\label{inset-docable-behavior}
DocAble distributes document accessibility work across workers. A worker can fail while producing a corrected document, and another worker may have to recover the job. Safe recovery therefore involves two related questions: when is a recovered artifact safe to publish as the result, and which worker currently has the right to complete the job? Separate models preserve the facts each question needs.

\textbf{Publishing a recovered result}

DocAble must preserve the ordering between publishing a recovered artifact and updating the database record that references it: the record may not name an artifact that has not yet been published. Reduced to its states, that concern is a small model.

\textbf{Recovery-ordering model · Behavioral}

\begin{itemize}
\tightlist
\item
  \textbf{Engineering question} --- In what order may recovery publish an artifact and update the record that refers to it?
\end{itemize}



\begin{itemize}
\tightlist
\item
  \textbf{Model} --- a transition system over the recovery lifecycle; the decisive fact is a transition that does \emph{not} exist.
\end{itemize}



\begin{itemize}
\tightlist
\item
  \textbf{Property} --- the database record may refer to the fallback artifact only after that artifact has been uploaded to persistent storage.
\end{itemize}



\begin{itemize}
\tightlist
\item
  \textbf{Quality attribute} --- data integrity; crash-safety.
\end{itemize}

% <!-- point: the-decisive-fact-is-an-absent-edge-between-staged-and-durable | The recovery obligation lives in the gap between local staging and durable publication: there is no direct FALLBACK_STAGED to DB_UPDATED edge, so the model excludes executions in which the record points at a not-yet-durable artifact. -->

A worker recovering a document job may have a corrected artifact available locally before it has been uploaded to persistent storage. These are different conditions: a local artifact may disappear if the worker fails, while a published artifact remains available to the service. The database record must not refer to the recovered artifact until publication succeeds. \autoref{recovery-ordering-model} reduces this obligation to three states. In \texttt{FALLBACK\_STAGED}, the artifact exists locally; in \texttt{FALLBACK\_UPLOADED}, it has been published to durable storage; and in \texttt{DB\_UPDATED}, the database refers to it. The permitted path passes through all three states. The forbidden shortcut from \texttt{FALLBACK\_STAGED} to \texttt{DB\_UPDATED} captures the failure the model must exclude: a worker could update the database and then crash before uploading the artifact, leaving the record pointing to a result that other workers cannot retrieve.

The model omits the storage API, retry machinery, and other recovery details because they do not determine this ordering requirement. To check the property, search for a path to \texttt{DB\_UPDATED} that does not pass through \texttt{FALLBACK\_UPLOADED}. In a larger lifecycle, checking only for the direct forbidden transition would miss an indirect route through other states. The model establishes publication order, not eventual completion: a worker could wait or retry indefinitely without violating the ordering constraint.

\begin{figure}[htbp]
\centering
\includegraphics[width=\linewidth]{book/latex/assets/recovery-ordering-model-70e2639fc3.pdf}
\caption{\emph{Recovery-ordering model.} A recovered artifact must be published to durable storage before the database may refer to it. The prohibited transition from \texttt{FALLBACK\_STAGED} directly to \texttt{DB\_UPDATED} illustrates the ordering violation; the property also excludes indirect paths that bypass \texttt{FALLBACK\_UPLOADED}.}
\label{recovery-ordering-model}
\end{figure}

% <!-- point: the-recovery-ordering-is-a-slice-of-a-larger-lifecycle | The recovery model keeps the states needed to check publication order; a fuller job lifecycle can support broader workflow questions when it includes the relevant transitions. -->

The three-state recovery model is a reduction for checking publication order. DocAble's fuller job lifecycle includes queued, leased, remediating, and completed or failed states. That model can address broader workflow questions when it includes the relevant transitions; the smaller recovery model keeps only what the ordering question needs.

\textbf{Distinguishing a current owner from a returning one}

DocAble leases each in-flight job to one worker. A lease carries an owner, a monotonically increasing generation, and a lifetime; when a lease looks abandoned it can be reclaimed and the job re-leased to another worker at the next generation.

\textbf{In-flight lease · Ownership}

\begin{itemize}
\tightlist
\item
  \textbf{Engineering question} --- Who currently owns this in-flight work, and can a returning stale owner disturb the current claim?
\end{itemize}



\begin{itemize}
\tightlist
\item
  \textbf{Model} --- a claim relating a worker to a job, stamped with a generation and a lifetime; reclaim advances the generation.
\end{itemize}



\begin{itemize}
\tightlist
\item
  \textbf{Property} --- a stale claim cannot clear or supersede the current one.
\end{itemize}



\begin{itemize}
\tightlist
\item
  \textbf{Quality attribute} --- correctness under concurrency; crash-safety.
\end{itemize}

% <!-- point: leased-iff-owner-is-the-models-definition | A job is leased exactly when it has an owner; the two are the same fact from two sides, so the model represents ownership as the presence of a current claim — this is what "leased" means, and it sets up the invariant that follows. -->

Start with what the model \emph{means}. A job is leased exactly when it has an owner; the two are the same fact seen from two sides, so the model represents ownership as the presence of a current claim rather than as a separate flag that could disagree with it. That is a definition, not a property that could fail. It sets up the invariant that can.

% <!-- point: the-epoch-generation-makes-a-stale-release-inapplicable | Suppose A holds generation 7, stalls, and is reclaimed; the job is re-leased to B at generation 8; a late release from A carrying generation 7 must not clear B's generation-8 claim, and the generation makes that decidable. | terms: soft-vs-hard-governance -->

The consequential invariant is about time and contention. Suppose worker A holds the claim at generation 7, appears to stall, and is reclaimed; the job is re-leased to worker B at generation 8. If A now wakes and tries to release using its generation-7 claim, that release must not clear B\textquotesingle s current generation-8 claim. The generation makes the distinction decidable: a release carrying an old generation names a claim that no longer owns, so it is inapplicable to the current one.

\autoref{in-flight-ownership-states} follows this sequence. Worker A initially holds generation 7. After its lease expires, the job becomes reclaimable, and worker B acquires generation 8. When A returns, its release request still names generation 7. The current claim belongs to B at generation 8, so the release must be rejected. The generation number distinguishes successive claims on the same job even when both workers believe they have authority to act. Without it, a delayed operation from A could clear B's valid claim. The model therefore preserves owner, generation, and lifetime while omitting the worker implementation and the machinery that stores the lease.

\begin{figure}[htbp]
\centering
\includegraphics[width=\linewidth]{book/latex/assets/in-flight-ownership-states-01b169485d.pdf}
\caption{\emph{Ownership under contention.} Worker A's generation-7 claim expires and worker B acquires generation 8. A later release from A cannot clear B's current claim. Recording the generation makes stale operations distinguishable from operations authorized by the current owner.}
\label{in-flight-ownership-states}
\end{figure}

The publication model separates a local artifact from a durably stored one. The ownership model separates a current claim from a superseded one. In each case, the useful addition is a distinction that changes whether a state change is acceptable. The models make those distinctions available for analysis; the later Alignment chapter examines the mechanisms that enforce them.
\end{bookbox}

% <!-- point: the-lease-model-omits-the-processing-topology | The model omits the processing topology: ownership, generation, lifetime, and reclamation are the consequential properties, and the implementation later changed how the job's work is partitioned without invalidating the model — the pointers to the code changed, not the model. -->

Notice what this model omits: the processing topology. Ownership, generation, lifetime, and reclamation matter to the question of who controls a job; how the job\textquotesingle s work is partitioned does not. DocAble\textquotesingle s implementation changed over time to improve performance while preserving those ownership semantics. We changed the pointers to the code, but not the model. The ownership properties could therefore remain invariants across those changes. The inset explains what an invariant requires and how a checker evaluates one.

% <!-- index-def: invariant -->

% <!-- box-family: inset -->

% <!-- inset-domain: FORMAL METHODS -->

% <!-- box-weight: callout -->

% mage-box-prefix: Inset —

\begin{bookbox}{inset}{What is an invariant?}
\label{inset-invariant}
An \textbf{invariant} is a property required to hold over a declared domain of states, structures, or observations. In a behavioral state machine, the invariant may be a predicate that must hold in every reachable state; in a structural model, it may require every observed dependency to belong to a declared set; in a measurement model, it may require a quantity to remain within a hard capacity bound. State both the property \emph{and} the domain the checker claims to establish it over.

Take a behavioral illustration. A job carries two fields at once --- a \texttt{status} of \texttt{free}, \texttt{leased}, or \texttt{done}, and a \texttt{lease\_owner} that names the holding worker or is empty. At the model level, leased-iff-owner was a definition; an implementation that represents the two sides as separately written fields turns it into a property that can fail, which is exactly why you write the invariant down. "A job is leased if and only if it has an owner" is a predicate over the pair: \texttt{(status~==~"leased")~==~(lease\_owner~is~not~None)}. A \texttt{free} status with a lingering \texttt{lease\_owner} is a stale lease; a \texttt{leased} status with no owner is a lost lease. Separating the two writes can make either inconsistent state reachable, and both break the predicate. The predicate becomes mechanically meaningful when a checker evaluates it over the declared reachable state space:

\begin{verbatim}
STATES = product({"free", "leased", "done"}, {None, "worker-1", "worker-2"})
def invariant(s):                                    # leased iff owned
    return (s.status == "leased") == (s.lease_owner is not None)
assert all(invariant(s) for s in reachable_states())  # check every reachable state
\end{verbatim}
\end{bookbox}

\section{One Representation, Several Questions}\label{heading:one-representation-several-questions}

% <!-- point: the-lease-answers-ownership-behavioral-and-measurement-questions-at-once | Read as ownership the lease says who controls a job; read as behavioral it traces a claimed-held-expired-reclaimed lifecycle; counted, it measures how much work is in flight — one representation, three questions. -->

A lease illustrates how one representation can answer several classes of question. Its owner identifies who currently controls a job. Its generation and lifetime distinguish current claims from stale ones and support reasoning about expiry and reclamation. Counting the current claims across jobs gives a measurement of how much work is in flight. These questions concern ownership, behavior, and capacity, yet they draw on the same represented facts.

% <!-- point: overlap-confirms-the-classes-sort-questions-not-data-structures | The overlap follows from the distinction introduced at the chapter's opening: model classes identify engineering questions, not mutually exclusive data structures — the observation the closing section builds on to connect the models. -->

The overlap follows from the distinction introduced at the opening of the chapter: model classes identify engineering questions, not mutually exclusive data structures. One representation can serve several questions when it preserves the facts each requires. \magevalue{sec:scenarios-and-the-joined-model} develops this idea further by connecting models through shared identities.

Ownership introduced permission implicitly: a stale claimant \emph{may not} clear a current lease. The next section examines decisions directly, including permissions and how to choose what should happen next.
